\documentclass[10pt,a4paper]{amsart}
\usepackage[margin=19mm]{geometry}
\usepackage{amsmath,amssymb,mathtools}
\usepackage[scr=rsfs,cal=euler]{mathalfa}
\usepackage{xfrac}
\usepackage[unicode,hidelinks,bookmarks=false]{hyperref}
\newcommand{\ie}{i.e.\ }
\newcommand{\cf}{cf.\ }
\newcommand{\resp}{respectively\ }
\newcommand{\Subsection}{\subsection}
\providecommand{\nobreakdash}{\nobreak\hbox{-}}
\setlength{\emergencystretch}{2em}
\multlinegap0pt
\usepackage[shortlabels]{enumitem}
\usepackage{amssymb}
\usepackage{cancel}
\usepackage[normalem]{ulem}
\usepackage{xcolor}
\usepackage{tikz}
\usepackage{mathtools}
\newtheorem{proposition}{Proposition}
\usepackage{cleveref}
\crefname{proposition}{Proposition}{Propositions}
\newcommand{\R}{\mathbb{R}}
\newcommand{\T}{\mathbb{T}}
\title{Approximate controllability of a bilinear wave equation and minimum time}
\author{Karine Beauchard, Thomas Perrin, Eugenio Pozzoli}
\date{Source: arXiv:2601.19544v1 (CC BY 4.0)}
\begin{document}
\maketitle

\noindent\emph{Source and licence.} Approximate controllability of a bilinear wave equation and minimum time. Version arXiv:2601.19544v1 (CC BY 4.0), distributed by arXiv under the Creative Commons Attribution 4.0 International licence, http://creativecommons.org/licenses/by/4.0/, as recorded in the arXiv licence metadata for this exact version and shown on the abstract page https://arxiv.org/abs/2601.19544v1. Full licence notice: \texttt{licenses/arxiv-2601.19544-CC-BY-4.0.txt}.\par\smallskip

\noindent\textit{Editorial scope.} Counted unit: the article's proposition Lem:M (source0.tex lines 633--649). The article's complete original proof of it is delivered with it (source0.tex lines 821--878, one \texttt{proof} environment of 58 lines, which the article places away from the statement). No mathematics is written, repaired or re-derived: the statement and the proof are the article's own lines, in the article's own notation, and no retained line is substituted or re-flowed.  Retained uncounted as the source-local context the proof needs: lines 115--121; lines 128--136; lines 720--724; lines 749--754.  Nothing was omitted from any retained span, so the package carries no omission record.  No retained line contains a citation, so the wrapper carries no bibliography and no bibliography entry.  No retained line contains a figure environment, an image-inclusion command or drawing code, so no image is delivered and the package declares no asset dependency.  Editorial formatting only, in addition: an AMS \texttt{amsart} preamble that restates the article's own declarations and macros used by the retained lines, namely the environments \texttt{proposition} on separate counters, together with the standard packages those lines need.  The retained environments print in this document as Proposition 1 in order of appearance, whereas the article prints the same items with the numbering of its own full document.  The complete native source of the article as distributed in the arXiv e-print for this exact version is bundled unchanged as \texttt{sources/193460/source0.tex}.
\begin{equation} \label{def:q_rho}
\rho = \left\lbrace \begin{array}{ll}
2 & \text{ if } d=1, \\
2^+ & \text{ if } d=2, \\
d & \text{ if } d \geq 3.
\end{array}\right.
\end{equation}

\begin{equation} \label{def:B}
A=\begin{pmatrix}
    0 & 1 \\ \Delta-1 & 0
\end{pmatrix},
\qquad
B=\begin{pmatrix}
    0 & 0 \\ 1 & 0
\end{pmatrix}.
\end{equation}

\begin{equation} \label{exp(phiB)}
e^{\phi B}=\begin{pmatrix}
    1 & 0 \\ \phi & 1 
\end{pmatrix}. 
\end{equation}

\begin{proposition} \label{Lem:strongCV}
Let $\rho$ be as in \eqref{def:q_rho} and $(m_1^{\tau},m_2^{\tau},m_3^{\tau})_{\tau \in [0,1]}$ be bounded a family of $W^{1,\infty} \times L^{\rho} \times L^{\infty}(\mathbb{T}^d)$ for which the associated operators $M_{\tau}$ (see \eqref{def:M}) strongly converge towards $M_0$ on $H^1 \times L^2(\T^d)$. We assume that $(e^{tM_0})_{t\in\R}$ is a group of bounded operators on $H^2 \times H^1(\T^d)$. Then the operators
$\exp(\tau A + M_{\tau})$ 
strongly converge towards 
$\exp(M_0)$ as $\tau \to 0$.
\end{proposition}

\begin{proposition} \label{Lem:M}
\begin{enumerate}[topsep=0cm, parsep=0cm, itemsep=0cm]
\item There exists $C>0$ such that, for every
$(m_1,m_2,m_3) \in  W^{1,\infty} \times L^{\rho} \times L^{\infty}(\mathbb{T}^d)$, the operator
\begin{equation} \label{def:M}
M:=\begin{pmatrix}
 m_{1} & 0 \\
 m_{2} & m_{3}
 \end{pmatrix}
\end{equation}
is bounded on $H^1 \times L^2(\T^d)$ and 
$\|M\| \leq C \big( \|m_1\|_{W^{1,\infty}}+\|m_{2}\|_{L^{\rho}}+\|m_{3}\|_{L^{\infty}} \big)$.
%
\item If $\phi \in L^{\rho}(\T^d)$ and $\nabla \phi \in L^{\widetilde{\rho}}(\T^d)$ then the operator
$\exp(\phi B)$ is bounded on $H^2 \times H^1(\T^d)$ (see \eqref{exp(phiB)}).
\end{enumerate}
\end{proposition} 

\begin{proof}
\noindent $\bullet$ We apply \Cref{Lem:strongCV} with $M_{\tau}=(\tau V+\phi) B$. By \Cref{Lem:M}, we have $\| M_{\tau}-M_0\| \leq C \|\tau V \|_{L^{\rho}} \to 0$ as $\tau \to 0$, i.e. $M_{\tau}-M_0$ converges to zero in $\mathcal{L}(H^1 \times L^2)$, thus $M_{\tau}$ strongly converges towards $M_0$ on $H^1 \times L^2(\T^d)$.
Moreover, by \Cref{Lem:M}, $(e^{t\phi B})_{t\in\mathbb{R}}$ is a group of bounded operators on $H^2 \times H^1(\T^d)$. Finally, by \Cref{Lem:strongCV}, $\exp(\tau A+M_{\tau})$ strongly converges towards $\exp(M_0)$ on $H^1 \times L^2$ when $\tau \to 0$, which gives the conclusion.   

\medskip

\noindent $\bullet$ We first prove that, for every $\gamma \in W^{1,\infty}(\T^d)$ and $t\in\R$,
\begin{equation} \label{formule1}
\exp(-\gamma B) \exp(t A_V) \exp( \gamma B) = \exp\big(t(A_V+\mathcal{M})\big) 
\quad \text{ where } \quad \mathcal{M}=\begin{pmatrix}
\gamma & 0 \\ -\gamma^2 & -\gamma
\end{pmatrix}
\end{equation}
(note that the coefficients of $VB+\mathcal{M}$ satisfy the assumptions of \Cref{Lem:M}).
Thanks to \eqref{def:B}, we obtain $\exp(\pm \gamma B)=I\pm\gamma B$, thus,
for $W_0=(w_1^0,w_2^0) \in H^1 \times L^2(\T^d)$, we have
$$\exp(-\gamma B) \exp(t A_V) \exp( \gamma B) W_0 =
\begin{pmatrix}
    w(t) \\ z(t)-\gamma w(t)
\end{pmatrix}
\quad \text{ where } \quad
\begin{pmatrix}
    w(t) \\ z(t)
\end{pmatrix}
= e^{t A_V} \begin{pmatrix}
    w_1^0 \\ w_2^0+\gamma w_1^0
\end{pmatrix}.$$
We have 
$$\frac{d}{dt} 
\begin{pmatrix}
    w \\ z-\gamma w
\end{pmatrix}
=
\begin{pmatrix}
    z \\ (\Delta+V)w-\gamma z
\end{pmatrix}
=\begin{pmatrix}
   \gamma w + (z-\gamma w) \\ (\Delta+V-\gamma^2)w-\gamma (z-\gamma w)
\end{pmatrix}
= (A_V+\mathcal{M}) \begin{pmatrix}
    w \\ z-\gamma w
\end{pmatrix},$$
and $(w,z-\gamma w)(0)=(w_1^0,w_2^0)$, thus
$$\begin{pmatrix}
    w(t) \\ z(t)-\gamma w(t)
\end{pmatrix}
= \exp \big(t(A_V+\mathcal{M}) \big) W_0$$
which proves \eqref{formule1}. Now, by applying \eqref{formule1} with $(t,\gamma) \leftarrow (\tau,\phi/\sqrt{\tau})$, we obtain
$$\exp \left(-\frac{\phi}{\sqrt{\tau}}B \right) \exp(\tau A_V) \exp \left( \frac{\phi}{\sqrt{\tau}} \right) 
= \exp(\tau A + M_{\tau}) 
\quad \text{ where } \quad
M_{\tau}=\begin{pmatrix}
\sqrt{\tau} \phi &  0 \\  \tau V  - \phi^2 & \sqrt{\tau} \phi
\end{pmatrix}.$$
We apply \Cref{Lem:strongCV}.
By \Cref{Lem:M},
$\|M_{\tau}-M_0\| \leq C ( \sqrt{\tau}\|\phi\|_{W^{1,\infty}} + \tau \|V\|_{L^{\rho}} ) \to 0$ as $\tau \to 0$, i.e. $M_{\tau}-M_0$ converges to zero in $\mathcal{L}(H^1 \times L^2)$ thus $M_{\tau}$ strongly converges towards $M_0$ on  $H^1 \times L^2(\T^d)$. Since $\phi^2 \in W^{1,\infty}(\T^d)$, by \Cref{Lem:M}, $(e^{tM_0})_{t\in\R}$ is a group of bounded operators on $H^2 \times H^1$. Finally, by \Cref{Lem:strongCV}, $\exp(\tau A+M_{\tau})$ strongly converges towards $\exp(M_0)$ on $H^1 \times L^2$ when $\tau \to 0$, which gives the conclusion.   
\end{proof}

\end{document}
